%% kapa: exam=2 round=35 year=2024 period=3 subject=compensation_law subject_ko=감정평가 및 보상법규 seq=2 number=2 points=30
%% source: Q-Net 감정평가사 기출문제 (cst00309) artlSeq=5213532 file=raw/1008/5213532/2024년 제35회 감정평가사 2차시험 문제지.pdf
%% conversion: kordoc 4.24.51 -> tools/gen2.py (auto); review=proofread-sonnet; flags=none
\begin{longquestion}{【문제 2】}
甲은 2023. 8. 23. 父로부터 A광역시 B구 소재의 토지(이하 ‘이 사건 토지’라 함)를 증여 받았고, 이 사건 토지에 관하여 증여 당시에는 2023. 1. 1.을 기준일로 하는 개별공시지가가 m\textsuperscript{2}당 2,200,000원으로 결정·고시되어 있었다. 甲은 이를 기초로 하여 산정한 증여세를 납부하고자 하였으나, 개별공시지가에 오류가 있음을 발견하여 「부동산 가격공시에 관한 법률」 제12조에 따른 개별공시지가 정정결정을 신청하였다. 그런데 B구의 구청장 乙은 甲의 정정결정신청에 대하여 정정불가 결정을 통지하였다. 한편 그 이후 乙은 이 사건 토지에 관하여 토지특성조사의 착오 등 지가산정에 잘못이 있다고 하여 B구 부동산가격공시위원회의 심의를 거쳐 위 개별공시지가를 m\textsuperscript{2}당 3,900,000원으로 정정하여 결정·고시하였다. 이에 관할 세무서장 丙은 이 사건 토지의 가액이 m\textsuperscript{2}당 3,900,000원이라고 보아 이를 기초로 증여재산의 가액을 산정하여 증여세부과처분을 하였다. 다음 물음에 답하시오. (단, 각 물음은 상호독립적임) (30점) \par
\medskip\noindent 물음 1) 甲이 乙의 정정불가 결정 통지를 대상으로 취소소송을 제기할 수 있는지를 설명하시오. (15점) \par
\medskip\noindent 물음 2) 甲은 乙의 개별공시지가 정정결정과 관련하여 i) 정정 사유가 있다고 하더라도 그 사유가 명백하여야만 비로소 정정할 수 있는데, 정정 사유가 명백하지 않음에도 불구하고 乙이 개별공시지가를 정정한 것은 위법하다고 주장하고 있다. 또한, ii) 설령 乙의 개별공시지가 정정결정이 타당하다고 하여도 이 사건 토지에 관하여 증여 당시 고시되어 있던 종전의 개별공시지가를 기초로 하지 아니한 丙의 증여세부과처분은 위법하다고 주장하고 있다. 甲의 주장이 타당한지에 관하여 각각 설명하시오. (15점) \par
\medskip
\begin{framedtext}
\textbf{참조 조문} \par
\textbf{<부동산 가격공시에 관한 법률>} \par
\textbf{제12조(개별공시지가의 정정)} 시장·군수 또는 구청장은 개별공시지가에 틀린 계산, 오기, 표준지 선정의 착오, 그 밖에 대통령령으로 정하는 명백한 오류가 있음을 발견한 때에는 지체 없이 이를 정정하여야 한다. \par
\textbf{<행정소송법>} \par
\textbf{제19조(취소소송의 대상)} 취소소송은 처분등을 대상으로 한다. 다만, 재결취소소송의 경우에는 재결 자체에 고유한 위법이 있음을 이유로 하는 경우에 한한다. \par
\end{framedtext}
\end{longquestion}
